\section{Инжекциите като средство за киберпрестъпления}
\label{sec:cybercrime}

Уводът поставя въпроса дали един механизъм, действащ по време на компилация, може
да изключи цял клас инжекции. Преди да се търси отговор, трябва да се покаже защо
въпросът има значение извън инженерната практика. Настоящата глава разглежда
инжекцията като средство за престъпления и като източник на отговорност за онези,
които обработват чужди данни. Първо се очертава мястото на инжекциите в
класациите на уязвимостите. След това се разглеждат три документирани случая и
правната рамка: международна, европейска и българска. Накрая се разглежда преходът
от наказване към превенция, който мотивира работата.

\subsection{Инжекцията в съвременния пейзаж на заплахите}
\label{sec:cybercrime-landscape}

\textbf{Инжекция} в широк смисъл е ситуация, при която недоверен вход бъде
изтълкуван от интерпретатор като част от команда, а не като данни. Формалното
определение, което работата използва, се въвежда в глава~\ref{sec:theory}.

Действащото издание на класацията на \textit{OWASP} за най-сериозните рискове при
уеб приложенията е \textit{OWASP Top 10:2025}~\cite{owasp_top10}. В него
инжекцията е на пето място (категория A05:2025), след като е слязла с две места
от трето място в предишното издание. Категорията обединява 37 класа слабости по
номенклатурата \textit{CWE} и според документа е с най-голям брой регистрирани
уязвимости (\textit{CVE}) от всички категории~\cite{owasp_a05_injection}.

Номенклатурата \textit{CWE} (\textit{Common Weakness Enumeration}) на \textit{MITRE}
описва видовете слабости, а не отделните уязвимости. Записът CWE-89 е озаглавен
„Improper Neutralization of Special Elements used in an SQL Command ('SQL
Injection')“~\cite{cwe89}. Записът CWE-943 е озаглавен „Improper Neutralization of
Special Elements in Data Query Logic“ и обхваща заявка към хранилище на данни,
например база от данни, без да се ограничава до SQL~\cite{cwe943}. Тази разлика в
обхвата е съществена за работата. В класацията \textit{2025 CWE Top 25 Most
Dangerous Software Weaknesses} CWE-89 е на второ място, с едно място нагоре спрямо
2024 година~\cite{cwe_top25_2025}. Това е най-новият списък, наличен на страницата
на проекта към 4 ноември 2026 година.

Тези източници не доказват, че инжекциите се използват за престъпления. Те
показват, че класът е широко разпространен, устойчив във времето и добре описан,
т.е. че защитата срещу него е въпрос на познати средства, а не на неизвестна заплаха.

\subsection{Документирани случаи}
\label{sec:cybercrime-cases}

За всеки случай има първичен или официален източник, който го свързва със SQL
инжекция. Цифри са посочени само там, където са прочетени в такъв източник.
Случаите илюстрират три правни положения: масово изнудване с откраднати данни,
наказателно преследване на извършител и административна санкция на администратор
на данни.

\subsubsection{MOVEit Transfer и CL0P (2023)}

Уязвимостта \textit{CVE-2023-34362} засяга приложението \textit{MOVEit Transfer} на
\textit{Progress Software}. Записът в Националната база данни за уязвимости на
\textit{NIST} я описва като SQL инжекция, която може да позволи на неудостоверен
нападател достъп до базата от данни, класифицира я като CWE-89 и ѝ дава оценка 9.8
по \textit{CVSS v3.1}~\cite{nvd_cve_2023_34362}. Съвместното предупреждение
\textit{AA23-158A} на \textit{FBI} и \textit{CISA} от 7 юни 2023 година приписва кампанията на
групата \textit{CL0P}. Според него нападателите са използвали уязвимостта, за да
инсталират уеб обвивката \textit{LEMURLOOT} и да изтеглят данни от базите на
приложението, а от 2021 година групата разчита главно на кражба на данни, а не на
криптиране. Използването е започнало на 27 май 2023 година според публични
източници, а на 2 юни 2023 година уязвимостта е добавена в каталога на активно
използваните~\cite{cisa_aa23_158a}.

\subsubsection{Heartland Payment Systems и Албърт Гонсалес (2009, 2010)}

На 17 август 2009 година прокуратурата на щата Ню Джърси съобщава за обвинителен
акт срещу трима души за пет корпоративни пробива, сред които \textit{Heartland
Payment Systems}, \textit{7-Eleven} и \textit{Hannaford}. Обвинителният акт описва
кражба на над 130 милиона номера на кредитни и дебитни карти от тези три компании и
определя атаките като „SQL-injection attack“. Съобщението напомня, че
обвинителният акт е само обвинение~\cite{doj_gonzalez_indictment_2009}. На 26 март
2010 година Министерството на правосъдието на САЩ съобщава, че Албърт Гонсалес е
осъден на 20 години и един ден лишаване от свобода по две обвинения в съзаговор
във връзка с достъпа до платежните мрежи на тези три компании. Присъдата е
едновременна с постановената предния ден в Бостън по друго
дело~\cite{doj_gonzalez_sentence_2010}. Съобщението за присъдата не използва
термина „SQL инжекция“. Приписването на метода се основава на обвинителния акт, а не
на съдебно установяване.

\subsubsection{TalkTalk и санкцията на ICO (2016)}

На 5 октомври 2016 година британската Служба на комисаря по информацията
(\textit{Information Commissioner's Office}, \textit{ICO}) налага на
\textit{TalkTalk} глоба от 400\,000 британски лири, определена от службата като
рекордна. Според нея атаката е осъществена от 15 до 21 октомври 2015 година чрез
три уязвими уеб страници в наследена инфраструктура. Нападателят е получил достъп до
личните данни на 156\,959 клиенти, като при 15\,656 от тях е имал достъп и до данни
за банкови сметки~\cite{ico_talktalk_2016}. Службата установява, че нападателят е
използвал SQL инжекция, че техниката е „well understood“, защити съществуват и
компанията е трябвало да знае за риска. Имало е и две незабелязани предупреждения:
успешна SQL инжекция на 17 юли 2015 година срещу същата уязвимост и втора атака
между 2 и 3 септември 2015 година. Нарушението е квалифицирано като неспазване на
седмия принцип на британския \textit{Data Protection Act 1998}, а отделно
наказателно разследване води Столичната полиция~\cite{ico_talktalk_2016}.
Санкцията е наложена по действалия тогава закон, а не по Общия регламент относно
защитата на данните. Тя е за липса на мерки срещу добре позната атака, а не за
самата атака.

Трите случая заемат три различни роли: обект на масово изнудване, извършител,
подведен под наказателна отговорност, и администратор, санкциониран за липса на
защита. Общото е, че слабостта е от клас, за който защитата е известна отдавна.
\textit{ICO} го казва за конкретния случай, а предупреждението на \textit{CISA} и
\textit{FBI} (раздел~\ref{sec:cybercrime-prevention}), което говори за широко огласена кампания срещу приложение за управляван трансфер на файлове, засегнала
хиляди организации, го казва за индустрията. То не назовава продукта~\cite{cisa_sqli_alert};
връзката с \textit{MOVEit} е извод на автора. Проблемът не е липсата на знание, а
това, че знанието се прилага ръчно, на всяко място поотделно.

\subsection{Правна рамка}
\label{sec:cybercrime-law}

Правното третиране на инжекцията се определя от резултата, а не от техниката.
Нападател, който чрез инжекция получи достъп до данни, ги промени или спре работата
на система, може да осъществи състава на престъпление по международни, европейски и
национални актове. Администраторът на данните носи самостоятелно задължение да
осигури защита. Прилагането на състав към конкретно деяние е въпрос на съдебна
преценка, затова връзката между инжекцията и даден състав по-долу е формулирана като
възможна правна квалификация. Съдебна практика не се разглежда.

\subsubsection{Конвенцията от Будапеща}

Конвенцията за престъпленията в кибернетичното пространство на Съвета на Европа
(\textit{ETS № 185}, Будапеща, 23 ноември 2001 година) е първият международен
договор в областта~\cite{budapest_convention}. Релевантни са три деяния от дял 1,
посветен на престъпленията срещу поверителността, цялостността и наличността на
компютърните данни и системи.\footnote{Заглавията на членовете на Конвенцията са
дадени в превод на автора, защото официален български текст не беше намерен.
Съдържанието е проверено по английския текст.} Член 2 (незаконен достъп,
\textit{illegal access}) задължава страните да установят като престъпление
умишления достъп до компютърна система или част от нея без право. Страната може да изисква деянието да е извършено
чрез нарушаване на мерки за сигурност или с нечестно намерение. Член 4 (намеса в
данни, \textit{data interference}) обхваща умишленото повреждане, изтриване,
влошаване, променяне или скриване на данни без право. Член 5 (намеса в система,
\textit{system interference}) обхваща умишленото сериозно възпрепятстване на
функционирането на система чрез въвеждане, предаване, повреждане, изтриване,
влошаване, променяне или скриване на данни. За случаи като \textit{Heartland} е
релевантен и член 8 (компютърна измама, \textit{computer-related fraud})~\cite{budapest_convention}.

\subsubsection{Директива 2013/40/ЕС}

Директива 2013/40/ЕС относно атаките срещу информационните системи задължава държавите членки
да предвидят като престъпления сходни деяния и замества Рамково решение 2005/222/ПВР~\cite{eu_2013_40}.
Подредбата на членовете е различна от тази на Конвенцията: в Директивата член 3
урежда незаконния достъп до информационни системи, член 4 незаконната намеса в
система, член 5 незаконната намеса в данни, а член 7 инструментите за извършване на
престъпления. При Конвенцията системата е в член 5, а данните в член 4.
Има и разлика по същество. Член 3 изисква деянието да е наказуемо, когато е
извършено чрез нарушаване на мярка за сигурност, поне в случаите, които не се
считат за маловажни, докато Конвенцията само дава на страните възможност да поставят
това условие. Инжекция, която заобикаля удостоверяването, може да бъде разглеждана
като преодоляване на мярка за сигурност, но това е въпрос на преценка по случая.
Член 9 изисква максимален срок на лишаване от свобода най-малко две години за
деянията по членове 3 до 7 и най-малко пет години за членове 4 и 5, когато са
извършени в рамките на престъпна организация, причиняват сериозни вреди или са
насочени срещу информационна система на критична инфраструктура~\cite{eu_2013_40}.

\subsubsection{Българският Наказателен кодекс}

Компютърните престъпления са уредени в глава девета „а“ „Компютърни
престъпления“, членове 319а до 319е на Наказателния кодекс~\cite{bg_nk}. Текстът е
прочетен на lex.bg към 4 ноември 2026 година и не е сверяван с „Държавен вестник“.
Посочват се само предметът на състава и размерът на лишаването от свобода. Глобите
не се разглеждат.

Член 319а, ал. 1 наказва неправомерния достъп до информационна система или части от
нея, в немаловажни случаи, с лишаване от свобода до шест години. За информация,
представляваща държавна или друга защитена от закон тайна, и за информационна
система, която е част от критична инфраструктура, ал. 4 предвижда от четири до осем
години. Член 319б, ал. 1 наказва неправомерното добавяне, копиране, използване,
променяне, пренасяне, изтриване, повреждане, влошаване, скриване или унищожаване на
компютърни данни или спирането на достъпа до тях, в немаловажни случаи, с лишаване от
свобода до шест години. При цел имотна облага ал. 3 предвижда от пет до седем години.
Член 319в предвижда до седем години, когато деянието по чл. 319б е по отношение на
данни, „които се дават по силата на закон, по електронен път или на електронен или
друг носител“. Член 319г наказва въвеждането на компютърен вирус или на друга
програма, предназначена за нарушаване на дейността на система или за узнаване,
заличаване, изменение или копиране на данни без разрешение (до седем години). Член
319д наказва създаването и разпространението на програми, пароли и кодове за достъп с
цел престъпление по чл. 171, ал. 3 или чл. 319а до 319г (до шест години). Член 319е
се отнася до нарушение на чл. 6, ал. 2, т. 5 от Закона за електронния документ и
електронните удостоверителни услуги~\cite{bg_nk}.

SQL инжекция, използвана за достъп до данни без право, може да осъществи състава на
чл. 319а, ал. 1. Когато чрез нея данни се копират, изменят или изтриват, възможна
квалификация е чл. 319б, ал. 1, а при цел имотна облага и ал. 3. И двата члена
се прилагат в немаловажни случаи. Това е авторов
анализ на текста на закона, а не извод от съдебна практика.

\subsubsection{Администраторът на данни: GDPR}

Българският текст на Регламент (ЕС) 2016/679 (\textit{GDPR}) не съдържа думата
„инжекция“. Член 32, параграф 1
задължава администратора и обработващия да прилагат подходящи технически и
организационни мерки за сигурност, съобразени с риска и с „достиженията на
техническия прогрес“, а буква г) включва процес на редовно изпитване и оценка на
ефективността им. Член 33 изисква уведомяване на надзорния орган без ненужно
забавяне и, когато е осъществимо, не по-късно от 72 часа след узнаването, освен ако
няма вероятност нарушението да породи риск за правата и свободите на лицата. Според
член 83, параграф 4, буква а) нарушенията на задълженията на администратора и
обработващия съгласно членове 8, 11, 25 до 39, 42 и 43, към които спадат членове 32
и 33, подлежат на административно наказание „глоба“ или „имуществена санкция“ в
размер до 10\,000\,000 евро или, в случай на предприятие, до 2\,\% от общия му
годишен световен оборот за предходната финансова година, която от двете суми е
по-висока~\cite{gdpr}. Връзката с инжекциите е тълкувателна: когато за клас уязвимости
съществува документирана и широко известна защита, нейната липса трудно може да се
обоснове с достиженията на техническия прогрес. Така разсъждава \textit{ICO} по
случая \textit{TalkTalk} по действалото тогава право~\cite{ico_talktalk_2016}. Член 33
показва и друго: правото предвижда уведомяване \textit{след} нарушението.

\subsubsection{NIS 2 и Законът за киберсигурност}

Директива (ЕС) 2022/2555 (\textit{NIS 2}) задължава съществените и важните субекти да
приемат подходящи и пропорционални мерки за управление на рисковете за мрежовите и
информационните системи (член 21, параграф 1). Сред минималните мерки е „сигурност
при придобиването на мрежови и информационни системи, разработване и поддръжка“,
включително, в английския текст, „vulnerability handling and disclosure“ (параграф 2,
буква д)~\cite{nis2}. Директивата се прилага към субекти от приложения I и II, които са най-малко
средни предприятия, а в определени случаи независимо от размера (член 2)~\cite{nis2}.
По данни на правна кантора и на проследяващ регистър в България тя е въведена чрез
изменение на Закона за киберсигурност, прието на 5 февруари 2026 година,
обнародвано на 13 февруари 2026 година и в сила от 17 февруари 2026
година~\cite{cms_bg_nis2,risi_bg_nis2}. Заглавната бележка на пречистения текст на
закона потвърждава само изменение, обнародвано в „Държавен вестник“, бр. 17 от 13-о
число на 2026 година~\cite{bg_zks}. Съдържанието на закона по членове не е
проверявано от първичен източник и не се цитира.

\subsection{От наказание към превенция}
\label{sec:cybercrime-prevention}

Разгледаните норми действат след събитието: наказателният закон санкционира
нападателя, регламентът изисква уведомяване след нарушението, а глобата на \textit{ICO}
е последваща. Нито една не намалява вероятността следващото приложение да съдържа
същата слабост.

На 25 март 2024 година \textit{CISA} и \textit{FBI} публикуват съвместно
предупреждение за проектиране с вградена сигурност (\textit{Secure by Design Alert})
„Eliminating SQL Injection Vulnerabilities in Software“. То определя SQL инжекциите
като устойчив клас дефекти, отбелязва, че индустрията знае как да ги елиминира от
десетилетия (като пример: \textit{MySQL} въвежда подготвените изрази през 2004
година), и призовава ръководителите на производителите да прегледат кода си и да
изключат целия клас дефекти от настоящите и бъдещите продукти~\cite{cisa_sqli_alert}.
Основната препоръка е параметризирани заявки с подготвени изрази, които отделят
SQL кода от данните на потребителя. Почистването на входа се оценява като крехко,
трудно за налагане в голям мащаб и често заобиколимо. Препоръчва се
параметризираните заявки да са стандартна практика, наложена в средата за
разработка, например чрез библиотеки, които правят сигурния път пътя по
подразбиране. Документът се позовава и на принципите за проектиране с вградена
сигурност: поемане на отговорност за резултатите за клиентите и изкореняване на
цели класове уязвимости, вместо отстраняване случай по случай~\cite{cisa_sqli_alert}.

Този аргумент е политически, но има техническа последица. Ако сигурният път трябва да
бъде единственият, въпросът е какво може една библиотека да направи невъзможно.
Предупреждението говори за SQL и за стойности на параметри. То не разглежда имената
на таблици и колони, структурата на заявката и хранилищата, които не са релационни.
Именно тук се намира изследователският въпрос на настоящата работа. Следващата глава
разглежда какво вече е предложено като защита и докъде достига то.
